\documentclass[aspectratio=169,10pt]{beamer}

\usepackage{fontspec}
\setsansfont{Arial}
\setmonofont{Menlo}[Scale=0.85]
\usepackage[english,vietnamese]{babel}
\usepackage{booktabs}
\usepackage{graphicx}
\usepackage{amsmath}
\usepackage{tikz}
\usetikzlibrary{arrows.meta,positioning}

% ============================================================
%  meta.tex — shared metadata for the report (main.tex) and slides (slides/slides.tex)
%  Edit here only; both documents \input this file.
% ============================================================
\newcommand{\metaUniversity}{Đại học Bách khoa Hà Nội}
\newcommand{\metaSchool}{Trường Công nghệ Thông tin và Truyền thông}
\newcommand{\metaCourse}{Seminar 1}
\newcommand{\metaTopicCode}{A1}
\newcommand{\metaTitle}{Dùng AI phát hiện mơ hồ, thiếu sót, mâu thuẫn trong tài liệu đặc tả yêu cầu}
\newcommand{\metaShortTitle}{A1 -- AI phát hiện lỗi đặc tả yêu cầu}
\newcommand{\metaLecturer}{TS. Nguyễn Thanh Hùng}
\newcommand{\metaLecturerUnit}{Khoa Khoa học Máy tính}
\newcommand{\metaGroup}{Nhóm 4}
\newcommand{\metaSemester}{Học kỳ 20252 --- Năm học 2025--2026}
\newcommand{\metaDate}{\today}

% Members table body (ID & Name & Role)
\newcommand{\metaMembers}{%
  20251175M & Nguyễn Tùng Dương & Nhóm trưởng \\
  20251177M & Trần Xuân Nam     & Thành viên \\
  20252636M & Phạm Gia Khánh    & Thành viên \\
}
% Short author line for slide footer / title
\newcommand{\metaAuthorsShort}{Nguyễn Tùng Dương, Trần Xuân Nam, Phạm Gia Khánh}

\graphicspath{{../pic/}{./}}
\usetheme{HUST}

\begin{document}

% --- Slide 1: Title ---
\begin{frame}[plain]
  \titlepage
  \note{[Dương] Kính chào thầy và các bạn. Hôm nay Nhóm 4 xin trình bày báo cáo Seminar 1, đề tài A1: Dùng AI phát hiện mơ hồ, thiếu sót và mâu thuẫn trong tài liệu đặc tả yêu cầu phần mềm.}
\end{frame}

% --- Slide 2: Agenda ---
\begin{frame}{Nội dung trình bày}
  \begin{enumerate}
    \item \textbf{Đặt vấn đề \& Bài toán A1:} Tầm quan trọng của chất lượng SRS và 4 câu hỏi nghiên cứu (RQ1--RQ4)
    \item \textbf{Phân loại lỗi \& Dữ liệu:} 15 tiểu loại lỗi, bộ dữ liệu PURE (107) + \texttt{seeded\_lms} (40)
    \item \textbf{Phương pháp thực nghiệm:} Bộ công cụ T0--T4, kỹ năng \texttt{SKILL.md}, chuẩn chuyên gia \& độ đo
    \item \textbf{Kết quả \& Phân tích lỗi:} \textit{(Cập nhật sau Phase 3--6)}
    \item \textbf{Hạn chế \& Kết luận:} Threats to validity, hướng phát triển và hỏi đáp
  \end{enumerate}
  \note{[Dương] Bài trình bày gồm 5 phần chính, đi từ bối cảnh bài toán đến thiết kế thực nghiệm, kết quả và thảo luận.}
\end{frame}

% --- Slide 3: Why it matters ---
\begin{frame}{Tại sao cần phát hiện sớm lỗi đặc tả yêu cầu (SRS)?}
  \begin{columns}[T]
    \begin{column}{0.54\textwidth}
      \begin{block}{Thực trạng đặc tả ngôn ngữ tự nhiên}
        \begin{itemize}
          \item Phần lớn tài liệu SRS viết bằng ngôn ngữ tự nhiên nên dễ chứa lỗi \textbf{mơ hồ}, \textbf{thiếu sót} và \textbf{mâu thuẫn}.
          \item Lỗi ở giai đoạn đầu lan truyền sang thiết kế, mã nguồn và kiểm thử với chi phí sửa chữa tăng theo cấp số nhân.
          \item Đọc duyệt thủ công tốn thời gian và dễ mắc lỗi \textit{subconscious disambiguation} (Berry et al., 2003).
        \end{itemize}
      \end{block}
    \end{column}
    \begin{column}{0.43\textwidth}
      \begin{block}{Giới hạn công cụ truyền thống}
        \begin{itemize}
          \item Công cụ dựa trên luật từ điển (QuARS, Smella) bắt tốt từ khóa bề mặt nhưng khó hiểu ngữ cảnh.
          \item Hạn chế khi kiểm tra mâu thuẫn logic chéo nhiều yêu cầu.
          \item \textbf{Cơ hội:} Tác tử AI kết hợp kỹ năng miền (\texttt{SKILL.md}) hỗ trợ kiểm định toàn diện hơn.
        \end{itemize}
      \end{block}
    \end{column}
  \end{columns}
  \note{[Dương] Nhấn mạnh chi phí sửa lỗi muộn và lý do các công cụ từ điển như QuARS/Smella chưa đủ để thay thế kiểm duyệt ngữ nghĩa.}
\end{frame}

% --- Slide 4: Problem A1 + RQ1-RQ4 ---
\begin{frame}{Bài toán A1 và các câu hỏi nghiên cứu}
  \begin{block}{Yêu cầu đề tài A1}
    Sử dụng AI phát hiện \textbf{mơ hồ (ambiguity)}, \textbf{thiếu sót (incompleteness)} và \textbf{mâu thuẫn (inconsistency)} trong SRS; đối chiếu với chuẩn chuyên gia trên dữ liệu công khai để đo \textbf{tỷ lệ phát hiện đúng (Recall)} và \textbf{tỷ lệ bỏ sót (Miss Rate)}.
  \end{block}
  \vspace{0.2cm}
  \begin{itemize}
    \item \textbf{RQ1 (Hiệu năng):} Các tác tử AI đạt Precision, Recall, Miss Rate và $F_1$ như thế nào so với chuẩn chuyên gia và công cụ nền tảng dựa trên luật?
    \item \textbf{RQ2 (Hiệu ứng kỹ năng):} Nạp \texttt{SKILL.md} (phân loại + checklist + schema) có cải thiện hiệu năng so với câu lệnh thuần (\textit{plain prompt})?
    \item \textbf{RQ3 (Theo loại lỗi):} Loại khiếm khuyết nào (mơ hồ / thiếu sót / mâu thuẫn) khó phát hiện nhất?
    \item \textbf{RQ4 (Độ ổn định):} Kết quả giữa các lượt chạy lặp có ổn định không?
  \end{itemize}
  \note{[Dương] Giới thiệu trọng tâm của A1 là một nghiên cứu thực nghiệm có đối chứng, trả lời 4 câu hỏi nghiên cứu RQ1 đến RQ4.}
\end{frame}

% --- Slide 5: Taxonomy + Real PURE examples ---
\begin{frame}{Bộ phân loại lỗi (15 tiểu loại) và ví dụ thực tế từ PURE}
  \small
  \begin{tabular}{p{2.35cm}p{3.8cm}p{6.05cm}}
    \toprule
    \textbf{Loại lỗi} & \textbf{Tiểu loại} & \textbf{Ví dụ thực tế từ PURE} \\
    \midrule
    \textbf{Ambiguity} & \texttt{vague\_term}, \texttt{referential}, \texttt{syntactic}, \texttt{coordination}, \texttt{weak\_phrase} & \textbf{\texttt{KP-019}:} \textit{``Sequence's length must not be more than a line (59 characters)''} / \textbf{\texttt{GJ-037}:} \textit{``easy to read error messages''}. \\
    \midrule
    \textbf{Incompleteness} & \texttt{tbd}, \texttt{missing\_actor}, \texttt{missing\_condition}, \texttt{missing\_error\_handling}, \texttt{missing\_value\_or\_unit}, \texttt{missing\_consequence} & \textbf{\texttt{PE-001}:} \textit{``The format for service information description is defined''} (bị động thiếu chủ thể định nghĩa). \\
    \midrule
    \textbf{Inconsistency} & \texttt{conflicting\_values}, \texttt{conflicting\_rules}, \texttt{terminology}, \texttt{duplicate} & \textbf{\texttt{PE-002} vs \texttt{PE-001}:} Lặp lại nguyên văn \textit{``The format for service information description is defined''}. \\
    \bottomrule
  \end{tabular}
  \vspace{0.15cm}
  {\footnotesize \textbf{Quy tắc chuẩn hóa:} Tối đa 1 phát hiện/\texttt{(req\_id, type)}; mâu thuẫn/trùng lặp neo vào \texttt{req\_id} xuất hiện sau và ghi ID trước vào \texttt{related\_ids}.}
  \note{[Nam] Trình bày bộ phân loại theo chuẩn ISO 29148 và Berry et al. (2003), minh họa bằng chính các câu lấy từ KeePass, Gamma J và Peering.}
\end{frame}

% --- Slide 6: Method overview TikZ ---
\begin{frame}{Tổng quan quy trình thực nghiệm}
  \centering
  \begin{tikzpicture}[
    node distance=0.45cm and 0.45cm,
    box/.style={rectangle, rounded corners=4pt, draw=hustred, fill=white, text width=2.05cm, align=center, minimum height=1.2cm, font=\scriptsize},
    freeze/.style={rectangle, rounded corners=4pt, draw=hustred, fill=hustred!12, text width=2.05cm, align=center, minimum height=1.2cm, font=\scriptsize\bfseries},
    arrow/.style={-{Stealth[length=2.2mm]}, thick, color=hustred}
  ]
    \node[box] (d) {\textbf{1. Dữ liệu}\\PURE (107)\\+ Seeded (40)};
    \node[box, right=of d] (e) {\textbf{2. Chuẩn GT}\\Claude 2 lượt\\+ $\kappa$ + phân xử};
    \node[freeze, right=of e] (f) {3. Đóng băng\\SHA-256\\(\texttt{HASHES.txt})};
    \node[box, right=of f] (t) {\textbf{4. Công cụ}\\T0 + T1--T4\\$\times$ 3 lượt lặp};
    \node[box, right=of t] (m) {\textbf{5. Đánh giá}\\\texttt{evaluate.py}\\P/R/Miss/$F_1$, CI};
    \draw[arrow] (d) -- (e);
    \draw[arrow] (e) -- (f);
    \draw[arrow] (f) -- (t);
    \draw[arrow] (t) -- (m);
  \end{tikzpicture}
  \vspace{0.35cm}
  \begin{block}{Đảm bảo tính khách quan}
    \begin{itemize}
      \item Chuyên gia (\textbf{Claude}) tách biệt họ mô hình với các công cụ đánh giá.
      \item Mọi công cụ chạy trong thư mục tạm chỉ đọc, hoàn toàn không thấy \texttt{ground\_truth/}.
      \item \texttt{evaluate.py} kiểm tra SHA-256 trong \texttt{HASHES.txt} trước khi tính độ đo.
    \end{itemize}
  \end{block}
  \note{[Nam] Giải thích cơ chế đóng băng SHA-256 trước khi chạy Phase 4 để tránh việc sửa đáp án sau khi thấy kết quả công cụ.}
\end{frame}

% --- Slide 7: Dataset ---
\begin{frame}{Bộ dữ liệu thực nghiệm (147 yêu cầu)}
  \begin{columns}[T]
    \begin{column}{0.55\textwidth}
      \small
      \begin{tabular}{llcp{2.3cm}}
        \toprule
        \textbf{Tài liệu} & \textbf{Nguồn} & \textbf{\#Req} & \textbf{Đặc điểm} \\
        \midrule
        \texttt{keepass} & PURE & 32 & Văn phong tự do, nhãn \texttt{REQ-xx} \\
        \texttt{peering} & PURE & 24 & Bị động, có câu trùng lặp \\
        \texttt{gammaj} & PURE & 51 & Nhiều từ mơ hồ giao diện \\
        \texttt{seeded\_lms} & Tự tạo & 40 & 10 sạch + 30 lỗi có kiểm soát \\
        \midrule
        \textbf{Tổng} & & \textbf{147} & \\
        \bottomrule
      \end{tabular}
    \end{column}
    \begin{column}{0.42\textwidth}
      \begin{block}{Thiết kế \texttt{seeded\_lms} (40 reqs)}
        \begin{itemize}
          \item \textbf{10 sạch:} Đầy đủ chủ thể, giá trị, đơn vị, xử lý lỗi.
          \item \textbf{10 Ambiguity:} 6 vague, 2 weak phrase, 1 referential, 1 \textit{and/or}.
          \item \textbf{10 Incompleteness:} 4 TBD, 3 missing consequence, 1 actor, 1 error, 1 value.
          \item \textbf{10 Inconsistency:} 5 giá trị, 4 quy tắc, 1 thuật ngữ (đối lập với câu sạch trước đó).
        \end{itemize}
      \end{block}
    \end{column}
  \end{columns}
  \note{[Nam] Nêu lý do kết hợp PURE (dữ liệu thực tế công khai) với \texttt{seeded\_lms} (dữ liệu mới chưa từng xuất hiện trong tập huấn luyện LLM).}
\end{frame}

% --- Slide 8: AI tools T0-T4 ---
\begin{frame}{Các cấu hình công cụ đối chứng (T0--T4)}
  \footnotesize
  \begin{tabular}{lllp{4.9cm}}
    \toprule
    \textbf{Mã} & \textbf{Công cụ / Mô hình} & \textbf{Chế độ} & \textbf{Cơ chế hoạt động} \\
    \midrule
    \textbf{T0} & \texttt{baseline\_rules.py} & Luật tất định & Từ điển mơ hồ, \textit{TBD}, bị động thiếu \texttt{by <Actor>}, trùng lặp Jaccard $\ge 0{,}80$. \\
    \textbf{T1 / T2} & \texttt{agy} (\texttt{gemini-4-hydrogen-high}) & Plain / +\texttt{SKILL.md} & Đối chứng tác động kỹ năng trên mô hình Gemini 4 Hydrogen. \\
    \textbf{T3 / T4} & \texttt{agy} (\texttt{gemini-4-neon-high}) & Plain / +\texttt{SKILL.md} & Đối chứng tác động kỹ năng trên mô hình Gemini 4 Neon. \\
    \bottomrule
  \end{tabular}
  \note{[Nam] Trình bày hai cặp mô hình Gemini (Hydrogen và Neon) tách biệt họ mô hình với chuyên gia Claude Opus, cùng chạy ở chế độ Plain và +SKILL.md.}
\end{frame}

% --- Slide 9: Expert ground-truth protocol ---
\begin{frame}{Giao thức xây dựng chuẩn chuyên gia (Claude Opus)}
  \begin{enumerate}
    \item \textbf{Pass 1:} Phiên Claude độc lập gắn nhãn toàn bộ 107 yêu cầu PURE $\rightarrow$ \texttt{expert\_pass1.json}.
    \item \textbf{Pass 2:} Tác tử Claude độc lập thứ hai (không thấy Pass 1 và \texttt{seed\_key.json}) gắn nhãn cả 4 tài liệu (147 yêu cầu) $\rightarrow$ \texttt{expert\_pass2.json}.
    \item \textbf{Đo tự nhất quán:} Script \texttt{check\_gt.py} tính \textbf{Cohen's $\kappa$} trên $107 \times 3 = 321$ cặp \texttt{(req\_id, type)} của PURE và kiểm tra Recall của Pass 2 trên \texttt{seeded\_lms}.
    \item \textbf{Phân xử:} Giải quyết mọi bất đồng giữa Pass 1 và Pass 2, ghi vào \texttt{adjudication\_log.md} $\rightarrow$ \texttt{final\_gt.json}.
    \item \textbf{Đóng băng:} Ghi mã băm SHA-256 vào \texttt{ground\_truth/HASHES.txt}.
  \end{enumerate}
  \note{[Khánh] Giải thích vai trò của Pass 2 trong việc kiểm định độ tin cậy trên tập cấy lỗi vì tác giả viết \texttt{seed\_defects.py} đã biết trước đáp án.}
\end{frame}

% --- Slide 10: Metrics ---
\begin{frame}{Luật khớp và các độ đo đánh giá}
  \begin{columns}[T]
    \begin{column}{0.48\textwidth}
      \begin{block}{Luật khớp}
        \begin{itemize}
          \item \textbf{Strict:} Khớp \texttt{(req\_id, type)}; với \texttt{inconsistency} cho phép khớp qua \texttt{related\_ids}. Khử trùng cặp đối xứng trước khi chấm.
          \item \textbf{Lenient:} Khớp nhị phân ở mức yêu cầu ($N = 147$: có lỗi hay sạch).
        \end{itemize}
      \end{block}
    \end{column}
    \begin{column}{0.49\textwidth}
      \begin{block}{Độ đo \& Ước lượng}
        \begin{itemize}
          \item $\text{Precision} = \frac{TP}{TP+FP}$, \quad $\text{Recall} = \frac{TP}{TP+FN}$
          \item $\text{Miss Rate} = 1 - \text{Recall}$, \quad $F_1 = \frac{2PR}{P+R}$
          \item \textbf{Bất định:} Bootstrap 95\% CI ($B = 1000$).
          \item \textbf{Ổn định:} Jaccard trung bình \& Fleiss' $\kappa$ qua 3 lượt.
        \end{itemize}
      \end{block}
    \end{column}
  \end{columns}
  \note{[Khánh] Trình bày công thức các độ đo, đặc biệt nhấn mạnh Recall (correct detection rate) và Miss Rate theo yêu cầu đề tài A1.}
\end{frame}

% --- Slides 11-16: Placeholders awaiting Phase 3-6 ---
\begin{frame}{Kết quả tổng hợp \& Độ tin cậy chuẩn (Chờ Phase 3--5)}
  \begin{block}{Trạng thái hiện tại}
    Công cụ nền tảng \textbf{T0} đã chạy xong trên cả 4 tài liệu (KeePass: 12, Peering: 21, Gamma J: 13, Seeded LMS: 19 phát hiện). Bảng số liệu đối chiếu P/R/Miss Rate/$F_1$ và biểu đồ \texttt{pic/*.png} sẽ được điền sau khi hoàn tất chuẩn chuyên gia Claude và chạy T1--T4.
  \end{block}
  \note{[Khánh] Cập nhật tiến độ: T0 đã có kết quả, đang chờ khôi phục quota Claude để chốt chuẩn GT và chạy T1-T4.}
\end{frame}

% --- Slide 17: Limitations ---
\begin{frame}{Hạn chế và mối đe dọa tính hợp lệ}
  \begin{itemize}
    \item \textbf{Chuyên gia AI thay cho chuyên gia người:} Đã giảm thiểu bằng 2 lượt Claude độc lập, $\kappa$ và tách biệt họ mô hình công cụ.
    \item \textbf{Rò rỉ dữ liệu PURE trong tiền huấn luyện:} Được kiểm soát bằng tài liệu mới hoàn toàn \texttt{seeded\_lms}.
    \item \textbf{Kích thước mẫu vừa phải (147 yêu cầu):} Lượng hóa bất định bằng Bootstrap 95\% CI ở mức yêu cầu.
    \item \textbf{Tính tự nhiên của lỗi cấy:} Các mẫu cấy trong \texttt{seeded\_lms} có cấu trúc rõ ràng hơn một số lỗi ngữ cảnh thực tế.
  \end{itemize}
  \note{[Khánh] Nêu thẳng thắn các hạn chế khoa học và biện pháp giảm thiểu đã áp dụng trong thiết kế.}
\end{frame}

% --- Slide 18: Conclusion & Q&A ---
\begin{frame}{Kết luận \& Hỏi đáp}
  \begin{block}{Tổng kết hiện trạng thực hiện}
    \begin{itemize}
      \item Đã hoàn thành toàn bộ hạ tầng dữ liệu (147 yêu cầu), kỹ năng \texttt{SKILL.md}, câu lệnh đối chứng, công cụ nền tảng \texttt{T0}, bộ sinh độ đo \& biểu đồ \texttt{evaluate.py} và khung báo cáo/slide HUST.
      \item Bước tiếp theo: Chạy chuẩn chuyên gia Claude (Phase 3), đóng băng SHA-256, chạy T1--T4 và hoàn thiện phân tích kết quả.
    \end{itemize}
  \end{block}
  \centering
  \vspace{0.3cm}
  {\Large\bfseries\color{hustred} Cảm ơn thầy và các bạn đã lắng nghe!}\\[0.15cm]
  {\small Mời thầy và các bạn đặt câu hỏi.}
  \note{[Dương] Cảm ơn thầy và các bạn; sẵn sàng trả lời câu hỏi.}
\end{frame}

% --- Slide 19: Live-demo backup slide ---
\begin{frame}{Phụ lục trình chiếu: Demo dự phòng (\texttt{SKILL.md} JSON)}
  \begin{columns}[T]
    \begin{column}{0.42\textwidth}
      \begin{block}{Đầu vào (\texttt{keepass.json})}
        \ttfamily\scriptsize
        \{\\
        \hspace*{1em}"doc\_id": "keepass",\\
        \hspace*{1em}"req\_id": "KP-019",\\
        \hspace*{1em}"text": "Sequence's length must not be more than a line (59 characters)"\\
        \}
      \end{block}
    \end{column}
    \begin{column}{0.55\textwidth}
      \begin{block}{Đầu ra JSON chuẩn theo \texttt{SKILL.md}}
        \ttfamily\scriptsize
        \{\\
        \hspace*{1em}"doc\_id": "keepass",\\
        \hspace*{1em}"findings": [\{\\
        \hspace*{2em}"req\_id": "KP-019",\\
        \hspace*{2em}"type": "ambiguity",\\
        \hspace*{2em}"subtype": "vague\_term",\\
        \hspace*{2em}"evidence": "more than a line",\\
        \hspace*{2em}"related\_ids": [],\\
        \hspace*{2em}"rationale": "Phrase 'more than a line' is subjective before adding (59 characters) and depends on display width."\\
        \hspace*{1em}\}]\\
        \}
      \end{block}
    \end{column}
  \end{columns}
  \note{[Dương] Slide dự phòng minh họa luồng 1 yêu cầu đầu vào chuyển thành JSON có cấu trúc theo đúng schema của SKILL.md.}
\end{frame}

\end{document}
